\documentclass[12pt]{article}
\usepackage[lmargin=1in,rmargin=1in,tmargin=1in,bmargin=1in]{geometry}

\usepackage{aryaman}

\setcounter{tocdepth}{2}

\title{On the Cohen--Macaulayness of abelian extensions of regular local rings}
\author{Aryaman Maithani}
\date{September 3, 2026}

\usepackage[backend=biber,style=alphabetic,doi=false,isbn=false,url=false,eprint=false,maxbibnames=5]{biblatex}
\addbibresource{talks.bib}

\begin{document}

\maketitle
% \tableofcontents

\section*{Introduction}

Notes I made for my talk at BIKES: \newline

\begin{blockquote}
	Let $R$ be a regular local ring with quotient field $K$. Let $L$ be a Galois extension of $K$ and $S$ the integral closure of $R$ in $L$. Assume that the Galois group is abelian. Is $S$ Cohen--Macaulay? Roberts showed the answer is affirmative when the order of the group is invertible in the residue field of $R$, and Koh gave a counterexample in mixed characteristic. Recently, in joint work with Singh and Sridhar, we give equicharacteristic counterexamples. We shall discuss all these results, being rather elementary constructions and proofs! 
\end{blockquote}

\section{Nonmodular}

	Paul Roberts \cite{Roberts:Abelian} showed the following:

	\begin{thm} \label{thm:roberts}
		Let $R$ be a \underline{regular local ring} with fraction field $K$. Let $L$ be a finite Galois extension of $K$ with an abelian Galois group. Assume moreover that the order of the Galois group is not divisible by the characteristic of the residue field of $R$. Let $S$ denote the integral closure of $R$ in $L$. Then~$S$ is a \underline{Cohen--Macaulay ring}.
	\end{thm}

	We observe some generalities. 
	The separability of $L/K$ implies that $S$ is a finite $R$-module. 
	Thus, the Cohen--Macaulayness of $S$ is equivalent to its $R$-freeness.\footnote{Use the Auslander--Buchsbaum formula, recalling that $R$ is regular.} 
	Let $G$ denote the Galois group of $L/K$. 
	The group $G$ acts on $L$ and $S$ while fixing $K$ and $R$, respectively. 
	Thus, $L$ is a $K[G]$-module and $S$ an $R[G]$-module. 
	Finiteness of $R \subseteq S$ gives us an isomorphism 
	$S \otimes_{R} K \cong L$, which is also one of $K[G]$-modules. 
	(The isomorphism is saying that every element of $L$ can be written as a fraction $s/r$ with $s \in S$ and $r \in R$.) 
	On the other hand, we also have $L \cong K[G]$ in view of the Normal Basis Theorem, giving us
	\begin{equation} \label{eq:L-KG}
		K[G] \cong S \otimes_{R} K. 
	\end{equation}

	We now outline the main ideas of the proof. 
	Roberts first reduces the problem to the case that $R$ is complete with an algebraically closed residue field $k$; we assume that this is now the case. 
	Writing $G$ as a sum of cyclic groups $\mathbb{Z}/d_{1} \oplus \cdots \oplus \mathbb{Z}/d_{k}$, we see that the group ring $k[G]$ is isomorphic to
	\begin{equation*} 
		k[G] = 
		\frac{k[x_{1}]}{(x_{1}^{d_{1}} - 1)}
		\times \cdots \times
		\frac{k[x_{k}]}{(x_{k}^{d_{k}} - 1)}.
	\end{equation*}
	The nonmodularity implies that each $d_{i}$ is coprime to $\chr(k)$. 
	As $k$ is algebraically closed, each factor above splits into linear factors, giving us that
	\begin{equation*} 
		k[G] \cong k_{1} \times \cdots k_{n}; 
		\quad \text{each $k_{i}$ isomorphic to $k$}.
	\end{equation*}
	As $R$ is complete, the above decomposition be lifted to $R[G]$, giving us
	\begin{equation*} 
		R[G] \cong R_{1} \times \cdots R_{n}; 
		\quad \text{each $R_{i}$ isomorphic to $R$}.
	\end{equation*}
	As $S$ is an $R[G]$-module, it must decompose as
	\begin{equation*} %\label{eq:S-decomp}
		S \cong S_{1} \times \cdots S_{n},
	\end{equation*}
	where each $S_{i}$ is an $R$-module. Moreover, the above decomposition is \emph{canonical}! 
	The module $S_{i}$ is given as the elements of $S$ annihilated by elements of 
	$R_{1} \times \cdots \times R_{n}$ whose $i$th component is zero. 

	In view of~\eqref{eq:L-KG}, we see that $K[G]$ can be decomposed in the following ways:
	\begin{align*} 
		K[G] \cong R[G] \otimes_{R} K &\cong 
		(R_{1} \otimes K) \times \cdots \times (R_{n} \otimes K)
		\cong K \times \cdots \times K, \\
		K[G] \cong S \otimes_{R} K &\cong 
		(S_{1} \otimes K) \times \cdots \times (S_{n} \otimes K).
	\end{align*} 
	As the above decompositions are canonical, we must have 
	$S_{i} \otimes K \cong K$ for each $i$. 
	However, $S$ is a reflexive $R$-module, and thus the same is true for $S_{i}$. 
	Thus, each $S_{i}$ is a reflexive $R$-module of rank one. 
	As $R$ is a UFD, this implies that each $S_{i}$ is isomorphic to $R$. 
	In turn,
	\begin{equation*} 
		S \cong S_{1} \times \cdots \times S_{n} \cong R^{n} 
		\quad
		\text{as $R$-modules.}
	\end{equation*}

	In fact, factoriality of $R$ was all that we used to conclude that $S$ is free. Thus, we obtain:

	\begin{por} \label{por:UFD}
		Let $R$ be a \underline{UFD} with fraction field $K$. Let $L$ be a finite Galois extension of $K$ with an abelian Galois group. Assume moreover that the order of the Galois group is not divisible by the characteristic of the residue field of $R$. Let $S$ denote the integral closure of $R$ in $L$. Then~$S$ is \underline{free over $R$}.
	\end{por}

	If we assume $G$ to only be nilpotent, the theorem fails \cite[Example~2]{Roberts:Abelian}. 

\section{Mixed characteristic}
	
	In the same paper, Roberts \cite[Example~1]{Roberts:Abelian} also showed that the hypothesis on the order cannot be dropped in Porism~\ref{por:UFD}: 
	\begin{ex}
		Consider the regular ring
		\begin{equation*} 
			R \coloneqq 
			\left(
			\frac{\widehat{\mathbb{Z}}_{2}[X, Y, U, V]}%
			{(Y^{2} - 4U - X^{2} V)}\right)_{(2, X, Y, U, V)}
			,
		\end{equation*}	
		where $\widehat{\mathbb{Z}}_{2}$ denote the ring of $2$-adic integers. 
		Then, one checks that $R$ is a UFD by showing that $2$ is prime and that $R[2^{-1}]$ is a UFD. 
		Letting $K$ denote $\Frac(R)$ and $L$ the Galois extension $K[\sqrt{V}]$, one checks that the integral closure $S$ is not free over $R$. 
		In fact, Roberts shows that no nonzero $S$-module can be free over $R$. 
	\end{ex}

	Six years later, Koh \cite[Example~2.4]{Koh:unramified} showed that the hypothesis on the order cannot be dropped in Theorem~\ref{thm:roberts} cannot be dropped either: 
	\begin{ex}
		Pick a regular local ring $R$ such that $\dim R \ge 3$ such that $\sqrt{-3}$ is part of a regular system of parameters of $R$.% 
		\footnote{In particular, $R$ is of mixed characteristic $3$.}
		Pick $x, y \in R$ such that $\sqrt{-3}, x, y$ is a part of a regular system of parameters. 
		% Let $a$ and $b$ be square free elements of $R$ such that
		% \begin{equation*} 
		% 	a \equiv x y^{4} \mod 27R \andd
		% 	b \equiv x^{4} y \mod 27R.
		% \end{equation*}
		Consider $a \coloneqq xy^4 + 27$ and $b \coloneqq x^4y + 27$, and define
		$\theta_1 \coloneqq \sqrt[3]{ab^2}$ and 
		$\theta_2 \coloneqq \sqrt[3]{a^2b}$.
		Let $S$ be the integral closure of $R$ in $L = K(\theta_1)$. 
		The extension $L/K$ is Galois of order $3$ but $S$ is not free over $R$. 
	\end{ex}

\section{Equicharacteristic}
	
	Forty years after \emph{that}, we have equicharacteristic counterexamples \cite[Theorem~2.2]{MSS:Abelian}. 
	Recall that if $G$ is a finite subgroup of $\GL_{n}(k)$, then $G$ acts on the polynomial ring $S \coloneqq k[x_1,\dots,x_n]$ via degree-preserving $k$-algebra automorphisms, where the action of $g \in G$ is given by
	\begin{equation*}
		\begin{bmatrix} x_1 & \hdots & x_n \end{bmatrix}^{\tr}
		\mapsto
		g \begin{bmatrix} x_1 & \hdots & x_n \end{bmatrix}^{\tr}.
	\end{equation*}
	Note the action of $G$ extends to $L \coloneqq \Frac(S)$ and we have 
	$L^{G} = \Frac(R) \eqqcolon K$, 
	the extension $L/K$ has Galois group $G$,
	and $S$ is the integral closure of $R$ in $L$. 

	\begin{ex} \label{ex:poly}
		Consider the action of $G = \left\langle \two{1}{1}{}{1} \right\rangle$ on $S \coloneqq \mathbb{F}_{p}[x, y]$. 
		The generator maps $x \mapsto x+y$ and $y \mapsto y$, using which we see that
		$R \coloneqq \mathbb{F}_{p}[x^p - x y^{p-1}, y] \subseteq S^{G}$. 
		One checks that $S$ is integral over $R$ and $[\Frac(S) : \Frac(R)]=p$ and that $R$ is a polynomial ring (hence normal). 
		Using this, one concludes that $S^G = R$, a polynomial ring. 
	\end{ex}

	We recall that an element $g \in G$ is a \deff{bireflection} if $\rank(g-1) \le 2$. 
	We have the following useful theorem of Kemper \cite[Theorem~2.1]{Kemper:CM}:
	\begin{thm}
		Let $K$ be a field of characteristic $p>0$, and $H$ a finite subgroup of~$\GL_n(K)$, with its natural action on the polynomial ring 
		$T \coloneqq K[x_1,\dots,x_n]$. 
		If $H$ is a $p$-group, and $T^H$ is Cohen-Macaulay, then $H$ is generated by bireflections.
	\end{thm}	

	We now construct our example in positive characteristic. 
	Consider the ring 
	\begin{equation*} 
		T \coloneqq \mathbb{F}_{p}[x_1, y_1, x_2, y_2, x_3, y_3]
	\end{equation*}
	with the natural action of the group $G \le \GL_{6}(\mathbb{F}_{p})$ consisting of matrices of the form
	\begin{equation*} 
		\phantom{\quad\text{ for } a,b,c \in \mathbb{F}_{p}.}
		\left[
		\begin{array}{cc|cc|cc}
		1 & a & & & &\\
		& 1 & & & &\\
		\hline
		& & 1 & b & &\\
		& & & 1 & &\\
		\hline
		& & & & 1 & c\\
		& & & & & 1
		\end{array}
		\right]
		\quad\text{ for } a,b,c \in \mathbb{F}_{p}.
	\end{equation*}
	In other words, we have taken three copies of Example~\ref{ex:poly}. 
	Within $G$, consider the subgroup generated by
	\begin{equation*} 
		\left[
		\begin{array}{cc|cc|cc}
		1 & 1 & & & &\\
		& 1 & & & &\\
		\hline
		& & 1 & 1 & &\\
		& & & 1 & &\\
		\hline
		& & & & 1 & 1\\
		& & & & & 1
		\end{array}
		\right],
	\end{equation*}	
	i.e., the diagonal copy of group from Example~\ref{ex:poly}. 
	Then, by Kemper's theorem, the ring of invariants $T^H$ is not Cohen--Macaulay, the only bireflection in $H$ being the identity. 
	On the other hand, the ring $T^G$ is a polynomial ring, stemming from Example~\ref{ex:poly}. 

	Thus, we have a obtained an extension
	\begin{equation*} 
		R \coloneqq T^G \subseteq T^H \eqqcolon S
	\end{equation*}
	of rings where the smaller ring is regular but the larger ring is not Cohen--Macaulay. 
	As $G$ was abelian, Galois theory tells us that $\Frac(S)/\Frac(R)$ is Galois with abelian Galois group $G/H \cong \mathbb{Z}/p \times \mathbb{Z}/p$. 
	One may also note that $S \subseteq T$ is Galois (at the field level) with Galois group $\mathbb{Z}/p$ with $S$ being a UFD but $T$ not $S$-free. 
	The rings above are constructed in the $\mathbb{N}$-graded realm, rather than the local one, but one may localise at the homogeneous maximal ideals to get the desired results. 

\printbibliography
\end{document}